\thispagestyle{plain}
\noindent
\begin{tikzpicture}
\node[fill=fnDark,text=white,rounded corners=2pt,inner sep=6pt,
      minimum width=\textwidth,anchor=west]
  {\large\bfseries FPGA-Neural --- Descrizione generale e caratteristiche};
\end{tikzpicture}

\vspace{6pt}
\noindent
{\small FPGA-Neural è un \textbf{acceleratore hardware parametrico per reti neurali}
feed-forward completamente contenuto nell'FPGA. Il calcolo (moltiplicazione,
accumulo, bias, attivazione, saturazione) avviene interamente on-chip in aritmetica
intera INT8/INT32; il sistema host fornisce solo configurazione, pesi, dati di
ingresso e controllo attraverso una semplice interfaccia SPI, senza mai far parte
del datapath computazionale. Un unico bitstream serve qualunque topologia fino al
massimo di build.}

\vspace{8pt}
\begin{multicols}{2}
{\color{fnDark}\large\bfseries Caratteristiche}\\[2pt]
{\footnotesize
\begin{itemize}[leftmargin=1.1em]
\item Datapath \textbf{INT8 $\times$ INT8 $\to$ INT16 $\to$ INT32}, accumulo a 32~bit
      con estensione di segno.
\item \textbf{Balanced binary adder tree} ($O(\log_2 \text{PARALLEL})$) al posto della
      riduzione lineare.
\item MAC parallelo configurabile: \code{PARALLEL} MAC hardware simultanei per neurone,
      mappati su DSP \code{MULT18X18D}.
\item Architettura completamente \textbf{parametrica}: \code{N\_INPUTS}, \code{N\_NEURONS},
      \code{PARALLEL}, \code{DATA\_WIDTH}, \code{ACC\_WIDTH}, \code{N\_LAYERS}.
\item \textbf{Larghezza di rete a runtime}: \code{n\_inputs\_real}/\code{n\_neurons\_real}
      per-layer, un solo bitstream per ogni topologia fino al massimo.
\item Attivazioni configurabili: \code{ACT\_RELU} (default) e \code{ACT\_NONE} (lineare
      con saturazione bilaterale), con saturazione INT8.
\item \textbf{Due tipi di rete}: classica multi-layer dense (\code{layer\_sequencer},
      buffer ping-pong) e \textbf{grafo arbitrario sparse} (\code{graph\_engine} +
      buffer di attivazione in block RAM \code{DP16KD}), selezionabili a runtime.
\item Sottosistema di \textbf{memoria dedicata}: interfaccia byte$\leftrightarrow$word,
      controller PSRAM parallelo asincrono con \textbf{page mode} (70~ns accesso
      casuale, 20~ns burst di pagina), 8~MB indirizzabili (23~bit).
\item Interfaccia host \textbf{SPI Mode 0} MSB-first, \code{SET\_NET\_TYPE}+dispatch, \code{STATUS.done}
      sticky/clear-on-read, \code{READ\_CONFIG} runtime.
\item \textbf{Sottosistema flash} boot/persistenza: accesso esclusivo della FPGA a una
      \code{W25Q128JV} SPI NOR (16~MB) via SPI master dedicato, copy engine
      flash$\leftrightarrow$PSRAM e catalogo a 16 slot con CRC32, 8 opcode host.
\item Verificato in \textbf{simulazione} (Icarus Verilog) e \textbf{sintesi reale}
      (Yosys + nextpnr-ecp5 + ecppack).
\end{itemize}}

\columnbreak

{\color{fnDark}\large\bfseries Applicazioni}\\[2pt]
{\footnotesize
\begin{itemize}[leftmargin=1.1em]
\item Inferenza a bassa latenza deterministica come periferica di
      SoC Linux, Raspberry-Pi-like, ESP32, microcontrollori.
\item Blocco hardware riusabile integrabile in progetti eterogenei
      (piattaforma, non singola rete).
\item Edge AI su reti dense compatte quantizzate INT8.
\item Off-loading del carico neurale dalla CPU host verso hardware
      dedicato con throughput prevedibile.
\end{itemize}}

\vspace{4pt}
{\color{fnDark}\large\bfseries Target \& toolchain}\\[2pt]
{\footnotesize
\begin{itemize}[leftmargin=1.1em]
\item FPGA: Lattice ECP5 \code{LFE5U-45F-8BG381C} ($-8$, CABGA381).
\item Sintesi: Yosys; place\&route: nextpnr-ecp5; bitstream: Project~Trellis
      (\code{ecppack}).
\item Simulazione: Icarus Verilog (\code{-g2012}).
\item PSRAM: ISSI \code{IS66WVE4M16EBLL-70BLI} (64\,Mb, 4M$\times$16).
\end{itemize}}
\end{multicols}

\vspace{2pt}
% --- tabella parametri chiave ---
\noindent
{\small\color{fnDark}\bfseries Parametri chiave (configurazione baseline caratterizzata)}
\vspace{2pt}

\noindent
\begin{tabularx}{\textwidth}{L{3.2cm}L{3.6cm}Y}
\toprule
\rowh \thd{Grandezza} & \thd{Valore} & \thd{Note} \\
\midrule
Precisione dati & INT8 (signed) & \code{DATA\_WIDTH}=8 \\
\rowa Accumulatore & INT32 (signed) & \code{ACC\_WIDTH}=32 \\
Ingressi / neuroni & 256 / 4 & baseline benchmark datapath \\
\rowa MAC simultanei & $2\ldots64$ & $=$\code{PARALLEL}$\times$\code{N\_NEURONS} \\
Attivazioni & ReLU, lineare & \code{ACT\_RELU} / \code{ACT\_NONE} \\
\rowa Fmax (P=2, datapath) & 87.88~MHz & benchmark datapath isolato \\
Fmax (P=2, sistema integrato) & 67.91~MHz & sistema completo incl. sottosistema flash, place\&route reale \\
Throughput MAC (P=16) & $\approx$3.34~G\,MAC/s & teorico, solo datapath \\
\rowa Memoria di lavoro & 8~MB PSRAM & bus parallelo 16-bit, 70~ns / 20~ns page mode \\
Spazio indirizzi & 23~bit (byte) & \code{ADDR\_WIDTH}=23 \\
\bottomrule
\end{tabularx}

\vspace{8pt}
\noindent
{\small\color{fnDark}\bfseries Diagramma a blocchi del sistema}
\begin{center}
\begin{tikzpicture}[node distance=6mm and 10mm,font=\footnotesize]
  \node[fnblockD,minimum width=26mm,minimum height=13mm] (host){HOST\\{\scriptsize configura / addestra / controlla}};
  \node[fnblockT,right=16mm of host,minimum width=52mm,minimum height=22mm] (eng){};
  \node[anchor=north,font=\footnotesize\bfseries,text=fnDark] at (eng.north){FPGA -- Neural Network Engine};
  \node[fnreg,fill=white] (spi) at ([yshift=-2mm]eng.center){\code{spi\_slave} + \code{spi\_engine}};
  \node[fnreg,fill=white,below=2.5mm of spi] (arb){\code{mem\_arbiter} + \code{layer\_sequencer}};
  \node[fnreg,fill=white,above=2.5mm of spi] (core){\code{neuron\_memory} $\to$ \code{neuron\_parallel} $\to$ \code{mac8}};
  \node[fnblock,right=16mm of eng,minimum width=24mm,minimum height=13mm] (ram){PSRAM 8\,MB\\{\scriptsize \code{psram\_controller}}};
  \draw[fnbus] (host) -- node[fnlbl,above]{SPI} (eng.west|-host);
  \draw[fnbus] (eng.east|-ram) -- node[fnlbl,above]{16-bit async} (ram);
\end{tikzpicture}
\end{center}
\begin{center}\footnotesize\itshape\color{fnGrey}
Il datapath neurale è interamente nell'FPGA; l'host non partecipa alle singole
operazioni MAC.\end{center}
